夜晚，柏宇和我走进了一家坐落于一家大厦二十一楼的爵士酒吧，从酒吧的窗边可以俯瞰整个城市的夜景，酒单上写着伍尔夫的小说《海浪》中的一句话，“我脚踏实地，亦随生命流动”。 柏宇照旧点了一杯加冰块的威士忌，而我则点了一杯用威士忌，梅斯卡尔，青柠混合的鸡尾酒，酒上来后，我尝了一口，入口比较凉爽，但是略微有些苦涩。

柏宇向来喜欢在酒吧见面，他认为一杯酒下肚，人们会自然而然卸去伪装，说出那些心里话。我原本不喜欢喝酒，我认为喝酒又苦涩又涨肚子，但最近我对喝酒的态度有所改观了，我认为酒本身并非重要，更多的是让人置身于一种退社会化的、轻松的氛围中，能够诱使人放下一切防备。

“相比于鸡尾酒，我更加喜欢纯饮，尤其喜欢威士忌，威士忌是通过农作物蒸馏出来的，是最纯粹的酒。”

酒吧里的爵士乐手刚刚演奏完了一首《So What》，这已经是我这个月第三次约柏宇出来，每次我都会问柏宇关于理佳的事情，与其说我想找到理佳，不如说我更关心的是理佳为什么会消失，我相信一个人即使再神秘，也总有办法找到她的踪迹，然而我却难以理解为什么一个人可以做出消失的事情，是为了同我断绝关系吗，这似乎有些自作多情了，难不成她竟然是鬼？突如其来的想法让我吓了一跳。

“最近的生活感觉有些无聊了，”，柏宇喝了一口酒，然后突然说道，“现在的工作轻松，也有足够闲暇时光，但总觉得缺一些东西。”

“会不会培养些兴趣爱好会比较好呢？”

“是啊，其实我也知道，有时候也在刻意去培养兴趣爱好，但又总觉得对什么都兴趣不大。年轻的时候喜欢踢足球，喜欢研究飞行模型，现在觉得赢也好，输也好，都没那么有所谓。工作后又很难结识新的人，总之，一切平淡如水。”

“这么看，好像也没什么办法。”

沉默有顷。

“说起来我好像还记得你高中时候好像在一个什么分享上面分享过一篇小说，好像是一篇......关于幻想的故事。”柏宇说。

“哦，你说的是《白夜》吧，陀思妥耶夫斯基的一篇小说，关于一个幻想家的故事。说起来有点讽刺，我当时还在嘲笑幻想家的可笑与卑微，现在想来......我似乎也变成了幻想家。”

我想起了有一节课上，老师要求我们每个人分享一篇自己喜欢的文学作品，教室里开着日光灯，窗外的天已经黑了，玻璃上映着我们自己的人影，我在讲台上分享了这个故事，当我讲完第四夜结束的时候停了下来，课堂下面的同学在窃窃私语，“怎么还会有这么卑微的人啊。”

“哈哈，你又在对号入座了。不过说实话，我觉得现代人可能或多或少都有一些，尤其是在互联网上真真假假的信息，总是惹的人心生涟漪。”

我又想起了一桩高中时候的往事，有一天我在学校里复习到很晚，突然感到憋闷难忍，感觉快要喘不过气来。我此刻为什么会在这里呢？学习的意义是什么？我死后能够在世界上留下什么东西呢？然后，莫名其妙地灵感迸发，整整写了一页的诗，可是又无人聆听。

我把这件事情告诉了柏宇。

“怪不得我觉得有段时间你总是不太高兴，”，柏宇沉吟了一下，然后说道，“虽然不太相同，但我好像也经历过相似的感受。我在高中时候一心想学习物理学，在大学时候也学习物理。可是大学生活却和我想象的大不一致，大学里很少有人真正将物理作为真正的事业，身边的同学要不就是因为选择学校而学习了物理，实际上要不就在浑浑噩噩中度过了大学时光，要不就寻找其他的路，老师们也未必认真对待，物理对他们而言只是工具罢了，是物理，或者数学，或者其他什么于他们而言似乎也不重要。于是我也放弃了在物理上面的钻研，也尝试同其他人一样寻找其他的道路，可找来找去也都没有太感兴趣的，我想，既然如此，那就回家吧。可总还是不甘心，这么多年的学习究竟是为什么呢？苦读多年却只能做一份人人都能做的工作，然而即使是这份工作也是来之不易的。”

柏宇不断用酒杯撞击着桌子，不协调的撞击声伴随着台上的音乐透露出局促不安。

我不知道该说些什么，这些话他之前从未说过，他在众人面前的形象一直是开朗与睿智的，然而背后却也蕴含了这么多的失望与踌躇。安慰一下他吗？这似乎显得太过轻浮，我没有经历他的生活，又怎么站到他的位置上想问题呢。批判一下当今的体制吗？这似乎也没什么意义，情绪的宣泄对当前的形势也不会有任何改变。

沉默良久，我终究还是说出了一句安慰柏宇的话，但是话刚说出口，我便觉得有些后悔，这句话轻飘飘地没什么重量，连我自己都觉得不知道该怎么回复，柏宇看了我一眼，没有说话，喝了一口酒，杯里的威士忌已经见底了，台上的爵士乐换了一首又一首，这是我们二人独有的沉默时光。窗外，城市的霓虹灯依旧亮着，夜晚的车流在一段静默后又多了起来，城市的男男女女流连于各个酒吧之中，台上的贝斯手低垂着头，手指在琴弦上缓慢地游走，音箱里传出闷钝的嗡鸣。

不过还有一件事，我希望能够问一下柏宇，尽管整个酒局已经接近尾声，在前面我们已经絮絮叨叨说了这么多自己的经历，然而到最后，我还是忍不住希望问一下。

“那个...... 柏宇，在你心里，你觉得理佳是个怎样的人呢？”我小心翼翼地问道。

“她啊，我觉得她既像天使，又像魔鬼。”

“为什么这么说呢？”

“她温柔、待人得体，和她在一起的时间会很开心，可是她很少说自己的想法，我总觉她心里有一团迷雾，我没有办法看到她心里在想什么，我甚至不知道她自己知不知道自己在想什么。之前你说她的突然消失，其实我也并不奇怪，可能是家里出了事，可能是想换一份工作了，即使是她想去北极看极光我也并不奇怪。”

柏宇说的话令我有所同感，在和理佳相处的过程中，我和理佳似乎只共享着表面的欢愉，至于在她的心里想着什么，我也是一概不知。之前我也同她探讨过这个问题，然而她却总是要么不回答，要不就将问题引到别的地方。

我赞同了柏宇，然后说了一下我的经历，“这么看来，似乎连你也不了解理佳，可是我总觉得有些奇怪，明明我们都在一个高中一个班，在同一个家乡里，共同上了这么多年的学，互相之间却也并不了解。可是我却又喜欢上了一个根本不了解的人，或者说，我以为因为我们有着共同的经历，所以我会了解她，原来到头来，我对她的了解也仅仅只比陌生人好一些。不过说起来，她难道对她的朋友也是如此吗，不透露任何的真实想法？”

“或许她会和她的朋友说，但或者在她的朋友面前，她会展示她的另一面，这其实我也并不清楚。”

我和柏宇互相搀扶坐上了出租车，我们两个人都有一些微醉，在最后我们又开始互相调侃取笑，在酒精作用下也逐渐口无遮拦了起来，我们回忆起以往的旧事，然而旧事说来说去也不过就那些，无非是一些过往的桃色八卦，或者是某些人的一些小秘密，然而新事却也没什么可以说的。不过最后，他还是提到了我的小说，

“说真的，我觉得你高中的小说写的太烂了，文采是有的，但是不知道在说什么，有点矫揉造作的感觉，不过现在可能写起来会好很多了。”

然后柏宇不说话了，再看过去，他已经睡着了。

我回到家后，时间已经很晚了，草草洗漱了一下准备睡觉，但总也睡不着，按理说一天的时光也足够劳累，但是大脑还是异常活跃，思绪很乱，但是又没办法不想，我翻了个身，面朝沙发靠背。靠背上有一块污渍，是很久以前打翻的咖啡，当时没有擦，后来就擦不掉了。我想明天请个假，但又懒得去请，在家中有什么事情可做呢？好像也没有什么，沉默似乎比忙碌更让人战栗。在半睡半醒之间，我站在悬崖边，一只脚踩空掉入悬崖，我突然惊醒，枕巾已经完全湿了，额头上全是汗珠。
